\section{Statistical Optimal Transport}
\label{sec-statistical-ot}
\subsection{Law of Large Numbers and Central Limit Theorem}
\label{sec-law-large-numbers-clt}
\label{sec-quantitative-clt}
\paragraph{Empirical laws.}
\begin{equation}\label{eq-empirical-law-alpha-n}
\hat\alpha_n \eqdef \frac1n\sum_{i=1}^n\delta_{X_i}.
\end{equation}
The ordinary law of large numbers says that empirical averages converge to expectations. In measure language this means that \(\hat\alpha_n\) converges weakly toward \(\alpha\), because testing \(\hat\alpha_n\) against a bounded continuous function \(\varphi\) gives the sample average \(n^{-1}\sum_i\varphi(X_i)\).

\begin{prop}[Empirical law of large numbers in \texorpdfstring{$\Wass_p$}{Wp}]\label{prop-empirical-lln-wasserstein}
Let \((\Xx,d)\) be a Polish metric space, let \(p\geq1\), and let \(\alpha\in\Pp_p(\Xx)\). Let \((X_i)_{i\geq1}\) be i.i.d. random variables with law \(\alpha\), and define \(\hat\alpha_n\) by~\eqref{eq-empirical-law-alpha-n}. Then
\[
\hat\alpha_n\rightharpoonup\alpha
\qandq
\Wass_p(\hat\alpha_n,\alpha)\longrightarrow0
\]
almost surely. Moreover, \(\EE\,\Wass_p(\hat\alpha_n,\alpha)^p\longrightarrow0.\)
\end{prop}
\begin{proof}
Fix a reference point \(x_0\in\Xx\), and write \(r(x)=d(x,x_0)\). Since \(\Xx\) is Polish, the weak topology on \(\Pp(\Xx)\) admits a countable convergence-determining class \((\varphi_k)_{k\geq1}\subset C_b(\Xx)\). For each fixed \(k\), the strong law of large numbers gives
\(\int \varphi_k\d\hat\alpha_n = \frac1n\sum_{i=1}^n\varphi_k(X_i) \longrightarrow \EE\varphi_k(X_1) = \int\varphi_k\d\alpha\)
almost surely. Intersecting the corresponding probability-one events over the countable set of indices gives, on a single event of probability one, convergence against every \(\varphi_k\). By the convergence-determining property, this is exactly \(\hat\alpha_n\rightharpoonup\alpha\).

The moment condition \(\alpha\in\Pp_p(\Xx)\) means \(\int r^p\d\alpha<+\infty\). Applying the strong law again to the integrable random variable \(r(X_1)^p\) gives
\(\int r^p\d\hat\alpha_n = \frac1n\sum_{i=1}^n r(X_i)^p \longrightarrow \int r^p\d\alpha\)
almost surely. Proposition~\ref{prop-wass-topology-polish} states that weak convergence plus convergence of the \(p\)-th moments is equivalent to \(\Wass_p\)-convergence on \(\Pp_p(\Xx)\). Hence \(\Wass_p(\hat\alpha_n,\alpha)\to0\) almost surely.

Set \(A_n=\int r^p\d\hat\alpha_n\) and \(M=\int r^p\d\alpha\). The triangle inequality through the Dirac mass \(\delta_{x_0}\), followed by \((a+b)^p\leq2^{p-1}(a^p+b^p)\), gives the deterministic bound
\(\Wass_p(\hat\alpha_n,\alpha)^p \leq 2^{p-1}\big(A_n+M\big).\)
The family \((A_n)_n\) is uniformly integrable. Indeed, by the de la Vall\'ee--Poussin criterion, choose a convex superlinear function \(\Psi\) such that \(\EE\Psi(r(X_1)^p)<+\infty\); Jensen's inequality gives
\(\EE\Psi(A_n) \leq \frac1n\sum_{i=1}^n \EE\Psi(r(X_i)^p) = \EE\Psi(r(X_1)^p).\)
Thus \((A_n+M)_n\), and hence \((\Wass_p(\hat\alpha_n,\alpha)^p)_n\), is uniformly integrable. Together with almost-sure convergence to zero, this implies \(\EE\Wass_p(\hat\alpha_n,\alpha)^p\to0\).
\end{proof}
\paragraph{Central-limit fluctuations.}
\begin{prop}[Berry--Esseen bound in $\Wass_1$]\label{prop-berry-esseen-w1}
Let $(X_i)_{i=1}^n$ be i.i.d. real random variables such that $\EE X_i=0$, $\EE X_i^2=1$ and $\EE |X_i|^3<+\infty$. If $\alpha_n$ is the law of $n^{-1/2}\sum_i X_i$ and $\gamma$ is the standard Gaussian law, then
\(\Wass_1(\alpha_n,\gamma) \leq \frac{C\,\EE |X_1|^3}{\sqrt n},\)
where $C$ is a universal constant.
\end{prop}
\begin{proof}
By Kantorovich--Rubinstein duality,
\[
\Wass_1(\alpha_n,\gamma)
=
\sup_{\Lip(h)\leq1}
\left|\EE h(S_n)-\EE h(G)\right|,
\qquad
S_n=n^{-1/2}\sum_i X_i,\ G\sim\gamma.
\]
For each such $h$, solve Stein's equation
\(f_h'(x)-x f_h(x)=h(x)-\EE h(G)\).
The solution satisfies \(\norm{f_h'}_\infty+\norm{f_h''}_\infty\leq C\) for a universal constant. Hence
\(\EE h(S_n)-\EE h(G)=\EE[f_h'(S_n)-S_nf_h(S_n)].\)
Write \(S_n^{(i)}=S_n-X_i/\sqrt n\). Independence and \(\EE X_i=0\) give
\[
\EE[S_nf_h(S_n)]
=
\frac1{\sqrt n}\sum_{i=1}^n
\EE\!\left[X_i\big(f_h(S_n^{(i)}+X_i/\sqrt n)-f_h(S_n^{(i)})\big)\right].
\]
Taylor's formula with integral remainder and \(\EE X_i^2=1\) imply
\[
\left|
\EE[S_nf_h(S_n)]-
\frac1n\sum_{i=1}^n\EE f_h'(S_n^{(i)})
\right|
\leq
\frac{C\,\EE|X_1|^3}{\sqrt n}.
\]
Moreover, the Lipschitz bound on \(f_h'\) yields
\[
\left|
\EE f_h'(S_n)-\frac1n\sum_{i=1}^n\EE f_h'(S_n^{(i)})
\right|
\leq
\frac{C\,\EE|X_1|}{\sqrt n}
\leq
\frac{C\,\EE|X_1|^3}{\sqrt n},
\]
where the last inequality uses \(\EE X_1^2=1\). Sharper constants and higher-order transport-distance refinements are available, but are omitted from the compact handout.
\end{proof}
\begin{prop}[Sharp symmetric lattice asymptotic]\label{prop-sharp-lattice-w1-clt}
Let \(X\) be centered, symmetric, of unit variance, and supported on a finite subset of \(a+h\ZZ\), where \(h>0\) is the maximal lattice span. If \(\alpha_n\) is the law of \(n^{-1/2}\sum_{i=1}^nX_i\) and \(\gamma=\mathcal N(0,1)\), then
\(\sqrt n\,\Wass_1(\alpha_n,\gamma)\longrightarrow \frac h4.\)
\end{prop}
\begin{proof}
Denote by \(F_n\) the CDF of \(\alpha_n\), and by \(\Phi\) and \(\varphi\) the standard Gaussian CDF and density. With the centered periodic sawtooth \(\psi(u)=\frac12-\{u\}\), the integrated lattice Edgeworth expansion gives
\[
F_n(x)-\Phi(x)
=
\frac h{\sqrt n}\,
\psi\!\left(\frac{\sqrt n\,x-na}{h}\right)\varphi(x)+r_n(x),
\qquad
\norm{r_n}_{L^1(\RR)}=o(n^{-1/2});
\]
see,, and. Vallender's one-dimensional identity states that \(\Wass_1(\alpha_n,\gamma)=\int_\RR|F_n-\Phi|\). Therefore \(\big||u+v|-|u|\big|\leq|v|\) reduces the claim to
\[
\int_\RR
\left|\psi\!\left(\frac{\sqrt n\,x-na}{h}\right)\right|
\varphi(x)\d x
\longrightarrow
\int_0^1|\psi(u)|\d u
=
\frac14.
\]
The convergence is periodic averaging: prove it first for compactly supported step functions, where it is a Riemann sum over periods, and conclude by \(L^1\)-approximation of \(\varphi\).
\end{proof}
\begin{prop}[Sharp symmetric density asymptotic]\label{prop-sharp-density-w1-clt}
Let \(X\) be centered, symmetric, of unit variance, with a density and moments of every order. Set \(\kappa_4=\EE[X^4]-3\) and \(H_3(x)=x^3-3x\). If \(\alpha_n\) is the law of \(n^{-1/2}\sum_{i=1}^nX_i\) and \(\gamma=\mathcal N(0,1)\), then
\(\Wass_1(\alpha_n,\gamma) = \frac{|\kappa_4|}{24n} \int_\RR |H_3(x)|\varphi(x)\d x +o(n^{-1}).\)
\end{prop}
\begin{proof}
The non-lattice Edgeworth expansion through order \(n^{-1}\) has an \(n^{-1/2}\) term proportional to the third cumulant, and \(n^{-1}\) terms proportional to the fourth cumulant and to the square of the third one. Symmetry makes the third cumulant vanish and gives
\(F_n(x)-\Phi(x) =-\frac{\kappa_4}{24n}H_3(x)\varphi(x)+r_n(x), \qquad \norm{r_n}_{L^1(\RR)}=o(n^{-1});\)
see and. Inserting this expansion into Vallender's identity and using \(\big||u+v|-|u|\big|\leq|v|\) proves the formula.
\end{proof}
\(\int_\RR |H_3(x)|\varphi(x)\d x =2\varphi(0)+8\varphi(\sqrt3) =\frac{2+8e^{-3/2}}{\sqrt{2\pi}},\)
\begin{equation}\label{eq-bernoulli-uniform-sharp-w1-clt}
\Wass_1(\alpha_n,\gamma)
\sim
\begin{cases}
\dfrac1{2\sqrt n},
& X\sim\frac12(\delta_{-1}+\delta_1),\\[2mm]
\dfrac{1+4e^{-3/2}}{10\sqrt{2\pi}}\dfrac1n,
& X\sim\operatorname{Unif}[-\sqrt3,\sqrt3].
\end{cases}
\end{equation}
\(\alpha_n=(D_{1/\sqrt n})_\sharp\alpha^{*n}, \qquad n\geq1,\)
\subsection{Sample Complexity}
\label{sec-sample-complexity}
\paragraph{Unregularized OT.}
\begin{equation}\label{eq-empirical-wasserstein-scale}
r_{n,p,d}
\eqdef
\begin{cases}
n^{-1/(2p)}, & d<2p,\\
n^{-1/(2p)}(\log(1+n))^{1/p}, & d=2p,\\
n^{-1/d}, & d>2p.
\end{cases}
\end{equation}
\begin{prop}[Empirical OT has intrinsic-dimension value rates]\label{prop-empirical-ot-rate}
Let \(p\geq1\), let \(\Xx\subset\RR^d\) be compact, and, for each \(\alpha\in\Pp(\Xx)\), let \(\widehat\alpha_n=n^{-1}\sum_{i=1}^n\de_{X_i}\) for i.i.d.\ samples \(X_i\sim\alpha\). Then
\begin{equation}\label{eq-empirical-wasserstein-moment-scale}
\sup_{\alpha\in\Pp(\Xx)}
\EE\!\left[\Wass_p(\widehat\alpha_n,\alpha)^p\right]^{1/p}
\leq
C_{\Xx,p,d}\,r_{n,p,d}.
\end{equation}
\[
\EE\!\left[
\abs{\Wass_p(\hat\alpha_n,\hat\beta_m)-\Wass_p(\alpha,\beta)}^p
\right]^{1/p}
\lesssim_{p,d}
r_{n,p,d}+r_{m,p,d}.
\]
\end{prop}
\begin{prop}[Dyadic partition bound for \texorpdfstring{$\Wass_1$}{W1}]\label{prop-dyadic-partition-w1}
Let \(\mathcal Q_j\) be nested dyadic partitions of \([0,1]^d\) into half-open cubes of side length \(2^{-j}\), with cubes on the outer face closed along the boundary of \([0,1]^d\). For all \(\alpha,\beta\in\Pp([0,1]^d)\) and every integer \(J\geq0\),
\(\Wass_1(\alpha,\beta) \leq \sqrt d \sum_{j=0}^{J-1}2^{-j} \sum_{Q\in\mathcal Q_{j+1}}\abs{\alpha(Q)-\beta(Q)} + \sqrt d\,2^{-J}.\)
\end{prop}
\begin{proof}
The proof is constructive: one decomposes the two measures into pieces which can be coupled at different spatial scales. Write \(\Omega=[0,1]^d\). We build positive measures \((\alpha_j,\beta_j)_{j=0}^J\) such that
\(\sum_{j=0}^J\alpha_j=\alpha, \qquad \sum_{j=0}^J\beta_j=\beta, \qquad \alpha_j(Q)=\beta_j(Q) \quad\text{for every }Q\in\mathcal Q_j.\)
Assuming such a decomposition for a moment, each pair \((\alpha_j,\beta_j)\) can be coupled inside the cells of \(\mathcal Q_j\). Indeed, on every cell \(Q\in\mathcal Q_j\) with positive mass, take the product coupling between the normalized restrictions of \(\alpha_j\) and \(\beta_j\) to \(Q\), multiplied by the common mass \(\alpha_j(Q)=\beta_j(Q)\). Summing over cells and over \(j\) gives a coupling of \(\alpha\) and \(\beta\). Since the diameter of a cell in \(\mathcal Q_j\) is at most \(\sqrt d\,2^{-j}\), this coupling has cost bounded by
\(\Wass_1(\alpha,\beta) \leq \sqrt d\sum_{j=0}^J2^{-j}\alpha_j(\Omega).\)

Start at the finest scale. For \(Q\in\mathcal Q_J\), define \(\alpha_J\) and \(\beta_J\) on \(Q\) by keeping the common part of the two masses:
\[
\alpha_J(A\cap Q)
=
\frac{\alpha(Q)\wedge\beta(Q)}{\alpha(Q)}\,\alpha(A\cap Q),
\qquad
\beta_J(A\cap Q)
=
\frac{\alpha(Q)\wedge\beta(Q)}{\beta(Q)}\,\beta(A\cap Q),
\]
with the convention that the corresponding term is zero when the denominator vanishes. For \(j=J-1,\ldots,1\), assume that the pieces at all finer levels \(k>j\) have been constructed and define the residuals
\(\alpha'_j=\alpha-\sum_{k=j+1}^J\alpha_k, \qquad \beta'_j=\beta-\sum_{k=j+1}^J\beta_k.\)
On each \(Q\in\mathcal Q_j\), set \(\alpha_j\) and \(\beta_j\) equal to the common part of \(\alpha'_j\) and \(\beta'_j\) inside \(Q\), exactly as above. Finally set \(\alpha_0=\alpha-\sum_{k=1}^J\alpha_k\) and \(\beta_0=\beta-\sum_{k=1}^J\beta_k\). Each step removes only common residual mass, so all measures are positive, and by construction \(\alpha_j(Q)=\beta_j(Q)\) on all cells \(Q\in\mathcal Q_j\).

For notational consistency, set \(\alpha'_0=\alpha_0\), \(\beta'_0=\beta_0\), \(\alpha'_J=\alpha\), and \(\beta'_J=\beta\). Let \(0\leq j<J\) and \(R\in\mathcal Q_j\). Since all pieces constructed at levels finer than \(j+1\) have equal mass on every child \(Q\in\mathcal Q_{j+1}\), the residual imbalance on \(Q\) just before the matching at level \(j+1\) is the original imbalance:
\(\alpha'_{j+1}(Q)-\beta'_{j+1}(Q) = \alpha(Q)-\beta(Q).\)
After removing the common part at level \(j+1\), the remaining source residual inside \(Q\) has mass
\(\alpha'_j(Q) = \alpha'_{j+1}(Q)-\alpha_{j+1}(Q) = \big(\alpha'_{j+1}(Q)-\beta'_{j+1}(Q)\big)_+.\)
Because \(\alpha_j\) is itself a common submeasure of the residual \(\alpha'_j\), the mass matched at level \(j\) inside \(R\) satisfies
\[
\alpha_j(R)=\beta_j(R)
\leq \alpha'_j(R)
=
\sum_{Q\subset R}
\big(\alpha'_{j+1}(Q)-\beta'_{j+1}(Q)\big)_+
\leq
\sum_{Q\subset R}
\abs{\alpha(Q)-\beta(Q)},
\]
where the sums are over \(Q\in\mathcal Q_{j+1}\). Summing over \(R\in\mathcal Q_j\) gives \(\alpha_j(\Omega) \leq \sum_{Q\in\mathcal Q_{j+1}}\abs{\alpha(Q)-\beta(Q)}.\)
The last level satisfies \(\alpha_J(\Omega)\leq1\). Substituting these bounds into the cost estimate proves the claim.
\end{proof}
\begin{proof}[Proof of Proposition~\ref{prop-empirical-ot-rate}]
General empirical Wasserstein moment estimates prove the one-sample statement for arbitrary \(p\); the two-sample estimate then follows from the reverse triangle inequality and Minkowski's inequality.
By the triangle inequality, \(\abs{\Wass_1(\hat\alpha_n,\hat\beta_m)-\Wass_1(\alpha,\beta)} \leq \Wass_1(\hat\alpha_n,\alpha)+\Wass_1(\hat\beta_m,\beta).\)
It is therefore enough to bound \(\EE\Wass_1(\hat\alpha_n,\alpha)\). Apply Proposition~\ref{prop-dyadic-partition-w1} to \((\hat\alpha_n,\alpha)\). If \(Q\in\mathcal Q_{j+1}\), then \(n\hat\alpha_n(Q)\) is a binomial random variable with parameter \(\alpha(Q)\), and hence
\(\EE\abs{\hat\alpha_n(Q)-\alpha(Q)} \leq \sqrt{\EE\abs{\hat\alpha_n(Q)-\alpha(Q)}^2} \leq \sqrt{\frac{\alpha(Q)}{n}}.\)
Summing over the \(2^{(j+1)d}\) cells and using the Cauchy--Schwarz inequality gives
\[
\sum_{Q\in\mathcal Q_{j+1}}
\EE\abs{\hat\alpha_n(Q)-\alpha(Q)}
\leq
\frac{1}{\sqrt n}
\sum_{Q\in\mathcal Q_{j+1}}\sqrt{\alpha(Q)}
\leq
\frac{2^{(j+1)d/2}}{\sqrt n}.
\]
Therefore, for every \(J\geq0\), \(\EE\Wass_1(\hat\alpha_n,\alpha) \lesssim_d 2^{-J} + \frac1{\sqrt n} \sum_{j=0}^{J-1}2^{j(d/2-1)}.\)
If \(d=1\), the sum is uniformly bounded and we choose \(J\) such that \(2^{-J}\lesssim n^{-1/2}\). If \(d=2\), the sum is \(J\), and choosing \(2^J\simeq n^{1/2}\) gives \((\log n)n^{-1/2}\). If \(d\geq3\), the sum is dominated by its last term; choosing \(2^J\simeq n^{1/d}\) gives
\(2^{-J} + n^{-1/2}2^{J(d/2-1)} \simeq n^{-1/d}.\)
This proves the displayed \(\Wass_1\) rates. The elementary dyadic estimate is not sharp for every law in the critical case \(d=2\). The proposition deliberately retains the distribution-free \((\log n)n^{-1/2}\) bound supplied by the dyadic proof.
\end{proof}
\paragraph{Lower bounds and minimax optimality.}
\begin{prop}[Minimax lower bound for bounded densities]\label{prop-minimax-lower-w1-density}
Let \(d\geq3\). Let \(\mathcal A\) be the class of probability measures on \([0,1]^d\) admitting a density \(\rho\) with respect to Lebesgue measure such that
\(\frac12\leq \rho(x)\leq \frac32 \qquad\text{for a.e. }x\in[0,1]^d.\)
For each \(n\geq1\), let \(\mathfrak E_n\) be the class of measurable estimators
\(\widetilde\alpha_n:([0,1]^d)^n\longrightarrow\Pp([0,1]^d).\)
There exists \(c_d>0\), depending only on \(d\), such that
\[
\inf_{\widetilde\alpha_n\in\mathfrak E_n}
\sup_{\alpha\in\mathcal A}
\EE_{(X_1,\ldots,X_n)\sim\alpha^{\otimes n}}
\left[
\Wass_1\big(\widetilde\alpha_n(X_1,\ldots,X_n),\alpha\big)
\right]
\geq
c_d\,n^{-1/d}.
\]
The expectation is over i.i.d.\ samples \(X_i\sim\alpha\). In particular, the high-dimensional rate of Proposition~\ref{prop-empirical-ot-rate} is minimax sharp for \(\Wass_1\), even when one allows arbitrary estimators rather than the empirical measure.
\end{prop}
\begin{proof}
We use a standard hypercube construction. Choose a smooth function \(\psi\) compactly supported in \((0,1)^d\), with \(\int\psi=0\), \(\norm{\psi}_\infty\leq1\), and whose positive and negative parts are separated in the first coordinate. There exists a compactly supported Lipschitz function \(\varphi\) such that \(\int\varphi\psi\d x>0\). We extend both functions by zero outside \((0,1)^d\), and absorb \(\Lip(\varphi)\) into the constants below.

Let \(m\) be an integer, \(h=1/m\), and partition \([0,1]^d\) into \(M=m^d\) cubes \(Q_k=z_k+h[0,1)^d\). Define the rescaled bumps
\[
\psi_k(x)=\psi\!\left(\frac{x-z_k}{h}\right)\mathbf 1_{Q_k}(x).
\]
For \(\theta=(\theta_1,\ldots,\theta_M)\in\{-1,1\}^M\) and a fixed small \(\eta>0\), set
\(\rho_\theta(x)=1+\eta\sum_{k=1}^M\theta_k\psi_k(x), \qquad \alpha_\theta=\rho_\theta\d x.\)
For \(\eta\) small enough, all these densities belong to \(\mathcal A\), and they integrate to one because every \(\psi_k\) has zero mean.

If \(\theta\) and \(\theta'\) differ on a set \(S\) of coordinates, the dual formula for \(\Wass_1\) gives an additive lower bound. Indeed, rescale \(\varphi\) as
\(\varphi_k(x)=h\,\varphi\!\left(\frac{x-z_k}{h}\right)\).
Each \(\varphi_k\) is supported in a fixed strict subcube of \(Q_k\). After dividing by this constant, Kantorovich--Rubinstein duality and the choice \(s_k=\operatorname{sign}(\theta_k-\theta'_k)\) on \(S\) give
\(\Wass_1(\alpha_\theta,\alpha_{\theta'}) \geq c\,\eta\,h^{d+1}\,\mathrm{Ham}(\theta,\theta'),\)
where \(\mathrm{Ham}\) denotes Hamming distance and \(c>0\) depends only on the chosen bump.

Neighboring vertices of the hypercube are statistically close. If \(\theta^{[k]}\) is obtained from \(\theta\) by flipping only the \(k\)-th sign, then the one-sample Kullback--Leibler divergence is bounded by
\(\KL(\alpha_\theta|\alpha_{\theta^{[k]}}) \leq C\eta^2 h^d,\)
because the two densities differ only on \(Q_k\), the perturbation has size \(O(\eta)\), and the densities are uniformly bounded away from zero. For \(n\) independent samples the divergence is multiplied by \(n\). Choose \(m\) so that \(m^d\simeq n\). Then \(n h^d\) is bounded, and, by taking \(\eta\) sufficiently small, Pinsker's inequality (Theorem~\ref{thm-pinsker}) gives a uniform bound
\(\TV\big(\alpha_\theta^{\otimes n},\alpha_{\theta^{[k]}}^{\otimes n}\big) \leq q<1.\)
Assouad's lemma therefore implies
\[
\inf_{\widetilde\alpha_n\in\mathfrak E_n}
\sup_{\theta\in\{-1,1\}^M}
\EE_{\alpha_\theta^{\otimes n}}
\left[
\Wass_1\big(\widetilde\alpha_n(X_1,\ldots,X_n),\alpha_\theta\big)
\right]
\geq
c'\eta\,M h^{d+1}(1-q)
=
c'\eta(1-q)h
\simeq
n^{-1/d}.
\]
Decreasing the constant handles the finitely many small values of \(n\), and the claim follows.
\end{proof}
\paragraph{Leveraging smoothness.}
\begin{prop}[Smoothed plug-in rates]\label{prop-smooth-plugin-w1-rate}
Let \(d\geq3\), \(0<s\leq1\), and let \(\alpha=\density{\alpha}\d x\), \(\beta=\density{\beta}\d x\) be probability measures on the unit-volume flat torus \(\TT^d\), with \(\max\{\norm{\density{\alpha}}_{H^s},\norm{\density{\beta}}_{H^s}\}\leq L\). Let \(\kappa\) be a smooth, symmetric probability density on \(\RR^d\), with finite second moment and rapidly decaying Fourier transform; a Gaussian is a typical choice. Denote by \(\kappa_h\) the periodization of \(h^{-d}\kappa(\cdot/h)\). For empirical measures \(\hat\alpha_n\) and \(\hat\beta_m\) formed from i.i.d. samples of \(\alpha\) and \(\beta\), define
\(\tilde\alpha_n \eqdef \hat\alpha_n\ast\kappa_{h_n}, \qquad \tilde\beta_m \eqdef \hat\beta_m\ast\kappa_{h_m}.\)
Then, for every \(0<h_n,h_m\leq1\),
\[
\EE\Wass_1(\tilde\alpha_n,\alpha)
\leq
C\left(h_n^{s+1}+n^{-1/2}h_n^{1-d/2}\right),
\]
and the analogous bound holds for \((\tilde\beta_m,\beta)\). Choosing \(h_n\simeq n^{-1/(d+2s)}\) and \(h_m\simeq m^{-1/(d+2s)}\) therefore gives
\[
\EE\abs{
\Wass_1(\tilde\alpha_n,\tilde\beta_m)
-
\Wass_1(\alpha,\beta)}
\leq
C\left(n^{-\frac{s+1}{d+2s}}+m^{-\frac{s+1}{d+2s}}\right).
\]
The constant depends only on \(d,s,L\) and \(\kappa\).
\end{prop}
\begin{proof}
It suffices to prove the one-sample bound. Set \(h=h_n\) and write \(\tilde\alpha_n=\density{\tilde\alpha_n}\d x\). For a zero-mean function \(w\) on the torus, set
\(\norm{w}_{\dot H^{-1}(\TT^d)}^2 \eqdef \sum_{k\in\ZZ^d\setminus\{0\}} \frac{\abs{\widehat w(k)}^2}{4\pi^2\norm{k}^2}.\)
Kantorovich--Rubinstein duality and Cauchy--Schwarz give
\(\Wass_1(u\d x,v\d x) \lesssim \norm{u-v}_{\dot H^{-1}(\TT^d)},\)
because a \(1\)-Lipschitz test function has \(L^2\)-gradient norm at most one on the unit-volume torus. Moreover, \(\EE\density{\tilde\alpha_n}=\density{\alpha}\ast\kappa_h\), so the two terms in the risk are a deterministic smoothing bias and a centered empirical fluctuation. Symmetry and the finite second moment of \(\kappa\) imply \(\abs{1-\widehat\kappa(\xi)}\lesssim\min\{1,\norm{\xi}^2\}\). Fourier estimates and Parseval's identity therefore give
\[
\begin{aligned}
\norm{\density{\alpha}\ast\kappa_h-\density{\alpha}}_{\dot H^{-1}}
&\lesssim h^{s+1}\norm{\density{\alpha}}_{H^s},\\
\EE\norm{\density{\tilde\alpha_n}-\density{\alpha}\ast\kappa_h}_{\dot H^{-1}}^2
&\lesssim \frac1n\sum_{k\neq0}\frac{\abs{\widehat\kappa(hk)}^2}{\norm{k}^2}
\lesssim \frac{h^{2-d}}n.
\end{aligned}
\]
The first estimate uses \(0<s\leq1\), and the last one uses \(d\geq3\) and the decay of \(\widehat\kappa\). The triangle and Jensen inequalities now give the one-sample bound. Equating its two terms yields \(h\simeq n^{-1/(d+2s)}\); the reverse triangle inequality for \(\Wass_1\) gives the two-sample claim.
\end{proof}
\(\tilde\alpha_{n,M} \eqdef \frac1{nM}\sum_{i=1}^n\sum_{j=1}^M \de_{X_i+h_n\xi_{i,j}}.\)
\[
\EE\Wass_1(\tilde\alpha_{n,M},\tilde\alpha_n)
\leq
h_n\,\EE\Wass_1\left(\frac1M\sum_{j=1}^M\de_{\xi_j},\Gaussian(0,\Id_d)\right)
\lesssim_d h_nM^{-1/d},
\]
\paragraph{MMD.}
MMD contains no transport optimization: its square is a combination of expectations of kernel evaluations, estimated empirically by sums of kernel values. Ordinary Monte-Carlo averaging therefore gives the dimension-free parametric rate below for bounded kernels.

\begin{prop}[MMD has a parametric value rate]\label{prop-mmd-sample-rate}
Let $k$ be a bounded positive definite kernel with RKHS $\RKHS_k$, and define
\(\MMD_k(\alpha,\beta) = \norm{\int k(x,\cdot)\d(\alpha-\beta)(x)}_{\RKHS_k}.\)
If $\hat\alpha_n$ and $\hat\beta_m$ are independent empirical measures, then
\[
\EE\abs{\MMD_k(\hat\alpha_n,\hat\beta_m)-\MMD_k(\alpha,\beta)}
\leq
\kappa\left(\frac1{\sqrt n}+\frac1{\sqrt m}\right)
\]
when $k(x,x)\leq\kappa^2$.
\end{prop}
\begin{proof}
Let $\Phi(x)=k(x,\cdot)$ be the feature map and $m_\alpha=\EE\Phi(X)$. The reverse triangle inequality for the RKHS norm gives
\(\abs{\MMD_k(\hat\alpha_n,\hat\beta_m)-\MMD_k(\alpha,\beta)} \leq \MMD_k(\hat\alpha_n,\alpha)+\MMD_k(\hat\beta_m,\beta).\)
Independence cancels the cross terms after taking the squared norm and expectation, giving
\(\EE\MMD_k(\hat\alpha_n,\alpha)^2 =\frac1n\EE\norm{\Phi(X)-m_\alpha}_{\RKHS_k}^2 =\frac{1}{n}\Big(\EE k(X,X)-\EE k(X,X')\Big).\)
The same estimate applies to $\hat\beta_m$, and Jensen's inequality together with $k(x,x)\leq\kappa^2$ gives the displayed bound.
\end{proof}
\paragraph{Entropic OT.}
\begin{prop}[Fixed-temperature Sinkhorn divergence has a parametric rate]\label{prop-sinkhorn-sample-rate}
Let \(c(x,y)=\norm{x-y}^2/2\), and assume that \(\alpha\) and \(\beta\) are \(\sigma^2\)-subgaussian probability measures on \(\RR^d\), in the sense that
\(\int e^{\norm{x}^2/(2d\sigma^2)}\d\alpha(x)\leq2, \qquad \int e^{\norm{y}^2/(2d\sigma^2)}\d\beta(y)\leq2.\)
Set
\[
q_d\eqdef\left\lceil\frac{5d}{2}\right\rceil+6,
\qquad
b_d\eqdef\left\lceil\frac{5d}{4}\right\rceil+3,
\qquad
\Lambda_{d,\sigma}(\epsilon)
\eqdef
\epsilon\left(1+\frac{\sigma^{q_d}}{\epsilon^{b_d}}\right).
\]
Let \(\hat\alpha_n\) and \(\hat\beta_m\) be independent empirical measures. Then, for every \(\epsilon>0\),
\[
\EE\abs{\bar\MK_\c^\epsilon(\hat\alpha_n,\hat\beta_m)-\bar\MK_\c^\epsilon(\alpha,\beta)}
\leq
C_d\,\Lambda_{d,\sigma}(\epsilon)
\left(\frac1{\sqrt n}+\frac1{\sqrt m}\right).
\]
\end{prop}
\begin{proof}

For probability measures \(\alpha_0,\beta_0\), write the dual objective, following~\eqref{eq-dual-sinkhorn-objective}, as
\[
\Dd_\epsilon^{\alpha_0,\beta_0}(u,v)
\eqdef
\int u\d\alpha_0+\int v\d\beta_0
-\epsilon\iint
\left(
\exp\left(\frac{u(x)+v(y)-c(x,y)}{\epsilon}\right)-1
\right)
\d\alpha_0(x)\d\beta_0(y).
\]
One can choose normalized optimal potentials for which the Sinkhorn equations hold pointwise on \(\RR^d\). Thus, at an optimal pair \((u,v)\),
\[
\int
\exp\left(\frac{u(x)+v(y)-c(x,y)}{\epsilon}\right)
\d\beta_0(y)
=1
\qquad\text{for all }x.
\]
Equivalently, the potentials are fixed points of the soft \(c\)-transform, but only the dual normalization is needed below.
Let \((u,v)\) and \((u_n,v_n)\) be optimal dual potentials for \((\alpha_0,\beta_0)\) and \((\alpha_n,\beta_0)\). By optimality,
\[
\Dd_\epsilon^{\alpha_0,\beta_0}(u_n,v_n)
\leq
\Dd_\epsilon^{\alpha_0,\beta_0}(u,v),
\qquad
\Dd_\epsilon^{\alpha_n,\beta_0}(u,v)
\leq
\Dd_\epsilon^{\alpha_n,\beta_0}(u_n,v_n).
\]
Hence the difference of the two optimal values is bounded by evaluating the two objectives at either pair of potentials. For the pair \((u,v)\) associated with \((\alpha_0,\beta_0)\), the preceding Sinkhorn equation cancels the exponential term in the change of the first marginal:
\(\Dd_\epsilon^{\alpha_0,\beta_0}(u,v) - \Dd_\epsilon^{\alpha_n,\beta_0}(u,v) = \int u\d(\alpha_0-\alpha_n).\)
The same identity holds for \((u_n,v_n)\), using its normalization with respect to \(\beta_0\).
\[
\abs{\MK_\c^\epsilon(\alpha_n,\beta_0)-\MK_\c^\epsilon(\alpha_0,\beta_0)}
\leq
2\sup_{w\in\mathcal F_{\epsilon,\sigma}}
\abs{\int w\d(\alpha_n-\alpha_0)},
\]
where \(\mathcal F_{\epsilon,\sigma}\) denotes the class of first-marginal entropic potentials arising from pairs of \(\sigma^2\)-subgaussian marginals.

For \(\epsilon=1\), one can show that the normalized quadratic-cost potentials belong to a polynomial-growth H\"older class. More precisely, after subtracting the quadratic part \(x\mapsto\norm{x}^2/2\), these potentials have derivatives up to order \(s_0=\lceil d/2\rceil+1\) whose local H\"older norms are bounded by polynomial functions of the subgaussian proxy. Denote by \(\mathcal F^s_{\sigma,L}\) such a class, where \(L\) is a random envelope controlled by the empirical subgaussian moment. Its empirical \(L_2\)-covering numbers satisfy a bound of the form
\(\log N(\delta,\mathcal F^s_{\sigma,L},L_2(\alpha_n)) \leq C_d\,L^{d/(2s)}\delta^{-d/s}(1+\sigma^{2d}).\)
This holds for \(s>s_0-1\), with \(\EE L\leq C\) after the normalization of the subgaussian moment. Since \(d/(2s_0)<1\), Dudley's entropy integral is finite, and the standard Rademacher symmetrization and chaining bound yields
\(\EE\sup_{w\in\mathcal F_{1,\sigma}} \abs{\int w\d(\hat\alpha_n-\alpha)} \leq C_d(1+\sigma^{q_d})\,n^{-1/2}.\)
This is exactly the step which avoids the exponential \(e^{D^2/\epsilon}\) factor: one controls the regularity class of the dual potentials polynomially instead of bounding their sup-norm through the worst value of the Gibbs kernel.

The general \(\epsilon\) case follows by scaling. If \(T_\epsilon(x)=x/\sqrt\epsilon\), \(\alpha_0^\epsilon=(T_\epsilon)_\sharp\alpha_0\), and \(\beta_0^\epsilon=(T_\epsilon)_\sharp\beta_0\), then
\(\MK_\c^\epsilon(\alpha_0,\beta_0) = \epsilon\,\MK_\c^1(\alpha_0^\epsilon,\beta_0^\epsilon).\)
The push-forward of a \(\sigma^2\)-subgaussian law by \(T_\epsilon\) is \((\sigma^2/\epsilon)\)-subgaussian, so the standard-deviation proxy \(\sigma\) is replaced by \(\sigma/\sqrt\epsilon\).
\[
\epsilon\,
C_d
\left(
1+
\left(\frac{\sigma}{\sqrt\epsilon}\right)^{q_d}
\right)n^{-1/2}
\leq
C_d\,\epsilon
\left(
1+\frac{\sigma^{q_d}}{\epsilon^{b_d}}
\right)n^{-1/2},
\]
where \(b_d=\lceil q_d/2\rceil=\lceil5d/4\rceil+3\) is the rounded exponent used in. Thus
\(\EE\abs{\MK_\c^\epsilon(\hat\alpha_n,\beta)-\MK_\c^\epsilon(\alpha,\beta)} \leq C_d\,\Lambda_{d,\sigma}(\epsilon)\,n^{-1/2}.\)
Changing both marginals is handled by the triangle inequality,
\[
\abs{\MK_\c^\epsilon(\hat\alpha_n,\hat\beta_m)-\MK_\c^\epsilon(\alpha,\beta)}
\leq
\abs{\MK_\c^\epsilon(\hat\alpha_n,\hat\beta_m)-\MK_\c^\epsilon(\alpha,\hat\beta_m)}
+
\abs{\MK_\c^\epsilon(\alpha,\hat\beta_m)-\MK_\c^\epsilon(\alpha,\beta)}.
\]
Empirical measures drawn from a subgaussian law have such a random proxy, with moments controlled by the population proxy. Integrating over this proxy gives the two-sample extension. The argument does not require the two empirical marginals to have the same sample size, so the two contributions scale as \(n^{-1/2}\) and \(m^{-1/2}\).

Finally, use the definition
\(\bar\MK_\c^\epsilon(\alpha,\beta) = \MK_\c^\epsilon(\alpha,\beta) -\frac12\MK_\c^\epsilon(\alpha,\alpha) -\frac12\MK_\c^\epsilon(\beta,\beta).\)
The cross term is controlled by the two-marginal estimate. For the self term, use the telescoping inequality
\[
\abs{\MK_\c^\epsilon(\hat\alpha_n,\hat\alpha_n)-\MK_\c^\epsilon(\alpha,\alpha)}
\leq
\abs{\MK_\c^\epsilon(\hat\alpha_n,\hat\alpha_n)-\MK_\c^\epsilon(\alpha,\hat\alpha_n)}
+
\abs{\MK_\c^\epsilon(\alpha,\hat\alpha_n)-\MK_\c^\epsilon(\alpha,\alpha)}.
\]
Each term is a one-marginal perturbation in the same pathwise sense: the other marginal may itself be empirical, but it is part of the common random subgaussian class. Thus the self term
\(\MK_\c^\epsilon(\hat\alpha_n,\hat\alpha_n)-\MK_\c^\epsilon(\alpha,\alpha)\)
is bounded at the same order as the cross term involving \(\alpha\), and similarly for \(\beta\). Summing these three contributions gives the displayed estimate, up to changing \(C_d\).
\end{proof}
\(\sigma=\frac{R}{\sqrt{2d\log 2}}=O(R).\)
\paragraph{Sample complexity of estimating OT maps.}
\[
u_{n,m}^\epsilon(x)
\eqdef
(\gD_{n,m}^\epsilon)^{\bar c,\epsilon}(x)
=
-\epsilon\log\sum_j b_j
\exp\!\left(\frac{\gD_{n,m,j}^\epsilon-c_2(x,Y_j)}{\epsilon}\right),
\qquad
T_{n,m}^\epsilon(x)\eqdef x-\nabla u_{n,m}^\epsilon(x).
\]
\[
\varphi_{n,m}^\epsilon(x)
\eqdef
\frac12\norm{x}^2-u_{n,m}^\epsilon(x)
=
\epsilon\log\sum_j b_j
\exp\!\left(
\frac{\dotp{x}{Y_j}+\gD_{n,m,j}^\epsilon-\frac12\norm{Y_j}^2}{\epsilon}
\right),
\qquad
T_{n,m}^\epsilon=\nabla\varphi_{n,m}^\epsilon.
\]
\begin{equation}\label{eq-empirical-entropic-map-barycentric}
T_{n,m}^\epsilon(x)
=
\sum_j \omega_j^\epsilon(x)Y_j,
\qquad
\omega_j^\epsilon(x)
=
\frac{
b_j\exp\!\left((\gD_{n,m,j}^\epsilon-c_2(x,Y_j))/\epsilon\right)
}{
\sum_k b_k\exp\!\left((\gD_{n,m,k}^\epsilon-c_2(x,Y_k))/\epsilon\right)
}.
\end{equation}
\(\omega_j^\epsilon(X_i)=\frac{\P_{i,j}^\epsilon}{a_i}, \qquad T_{n,m}^\epsilon(X_i) = \frac1{a_i}\sum_j \P_{i,j}^\epsilon Y_j.\)
Thus the extrapolated soft-transform map coincides on the empirical support with the barycentric projection~\eqref{eq-barycentric-projection} of the optimal entropic plan. At \(\epsilon=0\), the hard \(c\)-transform is generally nondifferentiable: away from Laguerre boundaries its gradient selects one target site, whereas on a boundary its subdifferential contains several sites.

\begin{prop}[Consistency of empirical barycentric maps]\label{prop-empirical-barycentric-map-consistency}
Let \(\Omega\subset\RR^d\) be compact, let \(\alpha,\beta\in\Pp_2(\Omega)\), and assume that \(\alpha\) is absolutely continuous. Let \(\phi\) be a Brenier potential from \(\alpha\) to \(\beta\) for the cost \(c_2(x,y)=\norm{x-y}^2/2\), and choose a bounded Borel selection \(T(x)\in\partial\phi(x)\) on \(\Omega\). At every differentiability point of \(\phi\), this selection necessarily equals \(\nabla\phi(x)\), so it is a representative of the unique Brenier map. Let
\(\hat\alpha_n=\frac1n\sum_{i=1}^n\delta_{X_i}, \qquad \hat\beta_m=\frac1m\sum_{j=1}^m\delta_{Y_j}\)
satisfy \(\hat\alpha_n\to\alpha\), \(\hat\beta_m\to\beta\) in \(\Wass_2\). Let \(\epsilon_{n,m}\geq0\) satisfy \(\epsilon_{n,m}\to0\) and \(\epsilon_{n,m}\log(\min\{n,m\})\to0\). For each \(n,m\), let \(\P_{n,m}^{\epsilon}\) be an optimal coupling between \(\hat\alpha_n\) and \(\hat\beta_m\), with \(\epsilon=\epsilon_{n,m}\): unregularized if \(\epsilon=0\), entropic if \(\epsilon>0\). Its barycentric projection, defined on the support of \(\hat\alpha_n\) by
\(\bar T_{n,m}^{\epsilon}(X_i) \eqdef n\sum_j \P_{i,j}^{\epsilon}Y_j,\)
satisfies \(\int \norm{\bar T_{n,m}^{\epsilon}(x)-T(x)}^2 \d\hat\alpha_n(x) = \frac1n\sum_i \norm{\bar T_{n,m}^{\epsilon}(X_i)-T(X_i)}^2 \longrightarrow0\)
as \(n,m\to+\infty\).
If \(\beta\) is also absolutely continuous, the analogous backward barycentric projections converge in \(L^2(\hat\beta_m)\) to the Brenier map \(S:\beta\to\alpha\).
\end{prop}
\begin{proof}
The assumptions imply \(\MK_{c_2}(\hat\alpha_n,\hat\beta_m)\to\MK_{c_2}(\alpha,\beta)\). If \(\epsilon=0\), \(\P_{n,m}^{\epsilon}\) is exactly optimal. If \(\epsilon>0\), compare its entropic objective with an unregularized optimal empirical plan \(Q_{n,m}\). Regarding \(Q_{n,m}\) as the joint law of two uniform discrete indices, its relative entropy with respect to their product law is their mutual information. Hence
\(\KLD(Q_{n,m}|\hat\alpha_n\otimes\hat\beta_m) \leq \min\{\log n,\log m\}.\)
It follows that
\(\int c_2\,\d \P_{n,m}^{\epsilon} \leq \MK_{c_2}(\hat\alpha_n,\hat\beta_m) +\epsilon\log(\min\{n,m\}).\)
Hence the transport costs of \(\P_{n,m}^{\epsilon}\) converge to the population optimum. By compactness, every weak limit point of the plans \(\P_{n,m}^{\epsilon}\) has marginals \(\alpha,\beta\) and optimal quadratic cost. Brenier uniqueness therefore forces the whole sequence to converge weakly to \((\Id,T)_\sharp\alpha\).

Disintegrate \(\P_{n,m}^{\epsilon}\) with respect to its first marginal and denote its barycentric projection by \(\bar T_{n,m}^{\epsilon}\). Jensen's inequality gives
\(\int \norm{\bar T_{n,m}^{\epsilon}(x)-T(x)}^2\d\hat\alpha_n(x) \leq \iint \norm{y-T(x)}^2\d \P_{n,m}^{\epsilon}(x,y).\)
At every differentiability point of \(\phi\), the subdifferential is the singleton \(\{\nabla\phi(x)\}\). Since the target is compact, the relevant subgradients are bounded, and the closed-graph property of \(\partial\phi\) implies continuity of the selected gradient at such a point. Convex functions are differentiable Lebesgue-almost everywhere; because \(\alpha\) is absolutely continuous, the bounded function \((x,y)\mapsto\norm{y-T(x)}^2\) is continuous \((\Id,T)_\sharp\alpha\)-almost everywhere. The portmanteau theorem for bounded functions whose discontinuity set has zero limiting mass gives
\(\iint \norm{y-T(x)}^2\d \P_{n,m}^{\epsilon}(x,y) \longrightarrow \int \norm{T(x)-T(x)}^2\d\alpha(x)=0.\)
The proof for the backward map is identical when the optimal plan is also induced by the Brenier map from \(\beta\) to \(\alpha\).
\end{proof}
\begin{prop}[Finite-sample rate for the entropic map]\label{prop-entropic-map-rate}
Assume that \(\alpha\) and \(\beta\) have densities on a common compact set \(\Omega\), both bounded above and with \(\density{\beta}\) bounded below away from zero. Let \(T=\nabla\phi\) be the quadratic Brenier map from \(\alpha\) to \(\beta\), and assume that \(\phi\in C^2(\Omega)\), its Legendre conjugate satisfies \(\phi^*\in C^{s+1}(\Omega)\) for some \(s>1\), and
\(\mu\Id_d\preceq\nabla^2\phi(x)\preceq L\Id_d \qquad (x\in\Omega)\)
for constants \(0<\mu\leq L\). Set \(\bar s\eqdef s\wedge3\), assume that \(\mathcal I_{\rm geo}(\alpha,\beta)<+\infty\) as in Proposition~\ref{prop-small-epsilon-expansion}, and draw \(n\) independent samples from each measure. If
\(\epsilon_n\simeq n^{-1/(d+\bar s+1)},\)
then the estimator \(T_{n,n}^{\epsilon_n}\) defined in~\eqref{eq-empirical-entropic-map-barycentric} satisfies
\[
\EE\norm{T_{n,n}^{\epsilon_n}-T}_{L^2(\alpha)}^2
\lesssim
\bigl(1+\mathcal I_{\rm geo}(\alpha,\beta)\bigr)
n^{-\frac{\bar s+1}{2(d+\bar s+1)}}\log n.
\]
\end{prop}
\paragraph{Sliced Wasserstein.}
For a direction \(\theta\in\Sphere^{d-1}\), let \(P_\theta(x)=\dotp{\theta}{x}\) be the one-dimensional projection and let \(\sigma\) denote the normalized surface measure on the sphere. For \(\alpha,\beta\in\Pp_p(\RR^d)\), the sliced Wasserstein distance is

\(\SW_p(\alpha,\beta)^p \eqdef \int_{\Sphere^{d-1}} \Wass_p\big((P_\theta)_\sharp\alpha,(P_\theta)_\sharp\beta\big)^p \d\sigma(\theta).\)
\begin{thm}[Dimension-free empirical rate for \texorpdfstring{$\SW_1$}{sliced Wasserstein-1}]\label{thm-sliced-sample-complexity}
Let \(\alpha,\beta\in\Pp(\RR^d)\) satisfy
\(\supp(\alpha)\cup\supp(\beta)\subset B(0,R)\), and let
\(\hat\alpha_n\) and \(\hat\beta_m\) be their empirical measures from \(n\) and \(m\) i.i.d. samples, respectively. Then
\[
\EE\abs{\SW_1(\hat\alpha_n,\hat\beta_m)-\SW_1(\alpha,\beta)}
\leq
R\left(\frac{1}{\sqrt n}+\frac{1}{\sqrt m}\right).
\]
\end{thm}
\begin{proof}
The triangle inequality for \(\SW_1\) gives \(\abs{\SW_1(\hat\alpha_n,\hat\beta_m)-\SW_1(\alpha,\beta)} \leq \SW_1(\hat\alpha_n,\alpha)+\SW_1(\hat\beta_m,\beta).\)
Fix \(\theta\in\Sphere^{d-1}\), set \(\alpha_\theta=(P_\theta)_\sharp\alpha\), and denote the cumulative functions of \(\alpha_\theta\) and its empirical measure by \(F_\theta\) and \(\hat F_{\theta,n}\). Both measures are supported on \([-R,R]\). The one-dimensional identity~\eqref{eq-w1-1d} and the Cauchy--Schwarz inequality give
\[
\EE\,\Wass_1\big((P_\theta)_\sharp\hat\alpha_n,\alpha_\theta\big)
=
\int_{-R}^R \EE\abs{\hat F_{\theta,n}(s)-F_\theta(s)}\d s
\leq
\int_{-R}^R \sqrt{\frac{F_\theta(s)(1-F_\theta(s))}{n}}\d s
\leq
\frac{R}{\sqrt n}.
\]
Indeed, \(n\hat F_{\theta,n}(s)\) is binomial with success probability \(F_\theta(s)\). Integrating this bound over \(\theta\) yields \(\EE\,\SW_1(\hat\alpha_n,\alpha)\leq R/\sqrt n\). The same argument for \(\beta\), followed by the preceding triangle inequality, proves the result.
\end{proof}
\subsection{Bias and Variance of OT}
\label{sec-bias-variance-ot}
\paragraph{Bias, variance and approximation errors.}
\(B_{n,m}^\epsilon \eqdef \EE \MK_c^\epsilon(\hat\alpha_n,\hat\beta_m) -\MK_c^\epsilon(\alpha,\beta)\)
\(Z_{n,m}^\epsilon \eqdef \MK_c^\epsilon(\hat\alpha_n,\hat\beta_m) -\EE \MK_c^\epsilon(\hat\alpha_n,\hat\beta_m).\)
\[
\MK_c^\epsilon(\hat\alpha_n,\hat\beta_m)-\MK_c^0(\alpha,\beta)
=
\underbrace{\MK_c^\epsilon(\alpha,\beta)-\MK_c^0(\alpha,\beta)}_{\text{regularization bias}}
+
\underbrace{B_{n,m}^\epsilon}_{\text{statistical bias}}
+
\underbrace{Z_{n,m}^\epsilon}_{\text{centered fluctuation}}.
\]
For the debiased Sinkhorn divergence, define \(\bar B_{n,m}^\epsilon\) and \(\bar Z_{n,m}^\epsilon\) by replacing \(\MK_c^\epsilon\) with \(\bar\MK_c^\epsilon\) above; the same decomposition then holds with bars. It is usually read after choosing a temperature \(\epsilon=\epsilon_n\).

\begin{prop}[Finite-space bias and CLT for OT]\label{prop-finite-ot-clt}
Let
\(\alpha=\sum_{i=1}^{N}a_i\delta_{x_i}, \qquad \beta=\sum_{j=1}^{M}b_j\delta_{y_j},\)
where all weights are positive and \(c\) is finite on the product support. Let \(\hat\alpha_n\) and \(\hat\beta_n\) be the independent empirical laws of \(n\) i.i.d. samples from \(\alpha\) and \(\beta\), respectively. Fix \(\epsilon\geq0\), with \(\MK_c^0=\MK_c\), and let \((f_\epsilon^\star,g_\epsilon^\star)\) be optimal dual potentials for \(\MK_c^\epsilon(\alpha,\beta)\). When \(\epsilon=0\), assume that this pair is unique up to the additive gauge; for \(\epsilon>0\), this uniqueness follows after fixing the gauge. Define
\[
\mathsf{v}_\epsilon
\eqdef
\sum_{i=1}^{N}a_i\bigl(f_\epsilon^\star(x_i)-\bar f_\epsilon\bigr)^2
+
\sum_{j=1}^{M}b_j\bigl(g_\epsilon^\star(y_j)-\bar g_\epsilon\bigr)^2,
\]
where \(\bar f_\epsilon=\sum_i a_i f_\epsilon^\star(x_i)\) and \(\bar g_\epsilon=\sum_j b_j g_\epsilon^\star(y_j)\). Then, as \(n\to+\infty\) with \(N,M\) fixed,
\[
\sqrt n\,B_{n,n}^\epsilon\longrightarrow0,
\qquad
\sqrt n\,Z_{n,n}^\epsilon\Longrightarrow\Gaussian(0,\mathsf{v}_\epsilon),
\qquad
n\,\operatorname{Var}\!\left[\MK_c^\epsilon(\hat\alpha_n,\hat\beta_n)\right]
\longrightarrow\mathsf{v}_\epsilon.
\]
Equivalently,
\[
\sqrt n\left(
\MK_c^\epsilon(\hat\alpha_n,\hat\beta_n)
-
\MK_c^\epsilon(\alpha,\beta)
\right)
\Longrightarrow\Gaussian(0,\mathsf{v}_\epsilon).
\]
For every fixed \(\epsilon>0\), one moreover has \(B_{n,n}^\epsilon=O(n^{-1})\).
\end{prop}
\begin{proof}
Writing \(\hat\a_n\) and \(\hat\b_n\) for the two empirical histograms, the independent multinomial central-limit theorems give
\(\sqrt n(\hat\a_n-\a,\hat\b_n-\b) \Longrightarrow (G_\a,G_\b),\)
where \(G_\a\) and \(G_\b\) are independent centered Gaussian vectors with covariance matrices \(\diag(\a)-\a\a^\top\) and \(\diag(\b)-\b\b^\top\). Proposition~\ref{prop-ot-first-variations-unregularized} gives the directional derivative for every \(\epsilon\geq0\); it is linear under the assumed uniqueness at \(\epsilon=0\), and automatically at positive temperature. Thus,
\(D\MK_c^\epsilon(\alpha,\beta)[h,k] = \sum_i f_\epsilon^\star(x_i)h_i + \sum_j g_\epsilon^\star(y_j)k_j.\)
The delta method therefore gives the Gaussian limit \(\sum_i f_\epsilon^\star(x_i)(G_\a)_i+\sum_j g_\epsilon^\star(y_j)(G_\b)_j\). Independence and the two multinomial covariance formulas give exactly \(\mathsf{v}_\epsilon\).

Set \(R\eqdef\max_{i,j}c(x_i,y_j)-\min_{i,j}c(x_i,y_j)\). Exact dual potentials may be replaced by their \(c\)-transforms, and entropic potentials by their soft \(c\)-transforms; in either case their oscillations are at most \(R\). Since differences of histograms have zero sum, integrating the corresponding subgradient bounds along segments in the two marginal simplices gives, for any histograms \((\a',\b')\) and their associated measures \((\alpha',\beta')\),
\(\abs{\MK_c^\epsilon(\alpha',\beta')-\MK_c^\epsilon(\alpha,\beta)} \leq \frac{R}{2}\bigl(\norm{\a'-\a}_1+\norm{\b'-\b}_1\bigr).\)
Because \(N,M\) are fixed, the fourth moments of \(\sqrt n\norm{\hat\a_n-\a}_1\) and \(\sqrt n\norm{\hat\b_n-\b}_1\) are uniformly bounded. The squared rescaled OT values are therefore uniformly integrable. The Gaussian limit then yields convergence of the means and second moments, which proves the assertions for \(B_{n,n}^\epsilon\), \(Z_{n,n}^\epsilon\), and the variance.

For fixed \(\epsilon>0\), the entropic value is twice continuously differentiable in a neighborhood of the positive pair \((\a,\b)\). On that neighborhood, a second-order Taylor expansion has a uniformly bounded remainder; its linear term has zero expectation and \(\EE(\norm{\hat\a_n-\a}^2+\norm{\hat\b_n-\b}^2)=O(n^{-1})\). The probability of leaving the neighborhood is exponentially small, and the preceding global Lipschitz bound controls its contribution. Hence \(B_{n,n}^\epsilon=O(n^{-1})\).
\end{proof}
\[
\sup_{(f,g)\in\mathcal D_c(\alpha,\beta)}
\left\{
\sum_i f(x_i)(G_\a)_i+
\sum_j g(y_j)(G_\b)_j
\right\},
\]
\paragraph{Fixed support versus a continuum.}
\[
\sum_{i=1}^{N}f_\epsilon^\star(x_i)(G_\a)_i
+
\sum_{j=1}^{M}g_\epsilon^\star(y_j)(G_\b)_j,
\qquad
\begin{cases}
\operatorname{Cov}(G_\a)=\diag(\a)-\a\a^\top,\\
\operatorname{Cov}(G_\b)=\diag(\b)-\b\b^\top.
\end{cases}
\]
\(\PP[\text{some source or target atom is unobserved}] \leq N e^{-n a_{\min}}+M e^{-n b_{\min}}.\)
\subsection{Sketching Sinkhorn in Linear Time}
\label{sec-sketching-sinkhorn}
\paragraph{PSD kernels and ordinary feature sketches.}
The RKHS/MMD section, Section~\ref{sec-rkhs-mmd}, already used positive semidefinite kernels to define Hilbertian discrepancies between measures. Recall that a symmetric kernel \(k:\Xx\times\Xx\to\RR\) is positive semidefinite if, for every finite family \((x_i)_{i=1}^n\), the Gram matrix \((k(x_i,x_j))_{i,j}\) is positive semidefinite.

\(k(x,y)=\int_Z \dotp{\phi(x,z)}{\phi(y,z)}_{\RR^p}\,d\xi(z), \qquad \phi:\Xx\times Z\to\RR^p,\)
\(\widetilde k_r(x,y) = \frac1r\sum_{\ell=1}^r \dotp{\phi(x,z_\ell)}{\phi(y,z_\ell)}_{\RR^p}.\)
\eql{\label{eq-rectangular-feature-sketch}
(\Phi_X)_{i,(\ell,q)}=r^{-1/2}\phi_q(x_i,z_\ell),
\qquad
(\Phi_Y)_{j,(\ell,q)}=r^{-1/2}\phi_q(y_j,z_\ell),
\qquad
\widetilde K=\Phi_X\Phi_Y^\top.
}
\(\widetilde K_{i,j} = \frac1r\sum_{\ell=1}^r \dotp{\phi(x_i,z_\ell)}{\phi(y_j,z_\ell)}_{\RR^p} = \widetilde k_r(x_i,y_j).\)
\begin{prop}[Entrywise accuracy of rectangular feature sketches]
\label{prop-sinkhorn-sketch-positive-guarantee}
Let \(z_1,\ldots,z_r\) be independent with law \(\xi\), and let \(K_{i,j}=k(x_i,y_j)\) and \(\widetilde K\) be the rectangular kernel matrix and sketch defined above, with \(\widetilde K\) given by~\eqref{eq-rectangular-feature-sketch}. Assume that, for some \(M<+\infty\),
\(\abs{\dotp{\phi(x_i,z)}{\phi(y_j,z)}_{\RR^p}}\leq M \qquad \text{for all }(i,j)\text{ and }\xi\text{-a.e. }z.\)
If \(K_{\min}\eqdef\min_{i,j}K_{i,j}>0\), then, for every \(0<\delta\leq1/2\),
\eql{\label{eq-rectangular-sketch-relative-concentration}
\PP\pa{
\max_{i,j}\abs{\frac{\widetilde K_{i,j}}{K_{i,j}}-1}>\delta
}
\leq
2nm\exp\pa{-\frac{rK_{\min}^2\delta^2}{2M^2}}.
}
On the complementary event, \(\widetilde K\) is entrywise positive and, for every \(\epsilon>0\), \(\max_{i,j} \abs{-\epsilon\log\widetilde K_{i,j}+\epsilon\log K_{i,j}} \leq 2\epsilon\delta.\)
\end{prop}
\begin{proof}
For each fixed \((i,j)\), the summands in the definition of \(\widetilde K_{i,j}\) have expectation \(K_{i,j}\) and lie in \([-M,M]\). Hoeffding's inequality therefore gives
\[
\PP\pa{
\abs{\widetilde K_{i,j}-K_{i,j}}>\delta K_{i,j}
}
\leq
2\exp\pa{-\frac{r\delta^2K_{i,j}^2}{2M^2}}
\leq
2\exp\pa{-\frac{r\delta^2K_{\min}^2}{2M^2}}.
\]
A union bound over the \(nm\) entries proves~\eqref{eq-rectangular-sketch-relative-concentration}. On its complementary event, write \(\widetilde K_{i,j}=K_{i,j}(1+s_{i,j})\), where \(\abs{s_{i,j}}\leq\delta\leq1/2\). Hence \(\widetilde K_{i,j}>0\), and \(\abs{\log(1+s_{i,j})}\leq2\abs{s_{i,j}}\) gives the logarithmic estimate.
\end{proof}
\paragraph{Application to Sinkhorn kernels.}
\(\alpha=\sum_{i=1}^n a_i\delta_{x_i}, \qquad \beta=\sum_{j=1}^m b_j\delta_{y_j}.\)
Specialize the preceding construction to the Gibbs kernel \(k_\epsilon(x,y)=e^{-c(x,y)/\epsilon}\), assumed to be PSD and to admit the displayed feature representation. Thus the rectangular matrix defined above is precisely the Sinkhorn kernel, \(K_{i,j}=k_\epsilon(x_i,y_j)\).

\(u=\frac{a}{Kv}, \qquad v=\frac{b}{K^\top u},\)
where divisions are componentwise. Replacing \(K\) by its rank-\(R\) sketch~\eqref{eq-rectangular-feature-sketch}, the two matrix-vector products are evaluated as

\(\widetilde K v=\Phi_X(\Phi_Y^\top v), \qquad \widetilde K^\top u=\Phi_Y(\Phi_X^\top u),\)
\begin{prop}[Accuracy and complexity of Gaussian Nystr\"om--Sinkhorn]
\label{prop-gaussian-nystrom-sinkhorn-complexity}
Let \(\alpha=\sum_{i=1}^n a_i\delta_{x_i}\) and \(\beta=\sum_{j=1}^m b_j\delta_{y_j}\) have positive weights, assume that all their support points lie in the ball \(B(0,D)\subset\RR^d\), and set \(N=n+m\). For the quadratic cost \(c(x,y)=\norm{x-y}^2\), fix \(\epsilon>0\) and \(\tau,\delta\in(0,1)\). There is a randomized Nystr\"om construction of an entrywise-positive rank-\(R\) approximation \(\widetilde K\) of the Gaussian kernel \(K_{i,j}=e^{-\norm{x_i-y_j}^2/\epsilon}\), followed by approximate Sinkhorn scaling and rounding, which returns a feasible coupling \(\widehat P\in\CouplingsD(a,b)\) and a value \(\widehat W\) such that, with probability at least \(1-\delta\),
\[
0\leq
\dotp{C}{\widehat P}
+\epsilon\KL(\widehat P\,|\,a\otimes b)
-\MK_c^\epsilon(\alpha,\beta)
\leq\tau,
\qquad
\abs{\widehat W-\MK_c^\epsilon(\alpha,\beta)}\leq\tau.
\]
The sketch rank can be chosen so that
\eql{\label{eq-gaussian-nystrom-rank}
R
\lesssim_d
\pa{
1+\frac{D^2}{\epsilon}
+\log\frac{N(1+D^2)}{\tau}
}^d
\log\frac{N}{\delta},
}
and the complete computation uses
\(\widetilde O\pa{ NR\pa{R+\frac{D^4}{\epsilon\tau}} } \quad\text{operations and}\quad O\pa{N(R+d)} \quad\text{memory}.\)
Here the constant hidden in \(\lesssim_d\) depends only on \(d\).
\end{prop}
\begin{proof}
Apply the Gaussian Nystr\"om approximation to the Gram matrix on the union of the \(N\) support points and retain its source--target block. Since \(\norm{x-y}\leq2D\), one has \(K_{\min}\geq e^{-4D^2/\epsilon}\). The adaptive stopping rule of Altschuler, Bach, Rudi and Niles-Weed resolves the kernel below this scale, hence produces \(\widetilde K>0\) with the same logarithmic control singled out in Proposition~\ref{prop-sinkhorn-sketch-positive-guarantee}. Their Gaussian effective-dimension estimate gives~\eqref{eq-gaussian-nystrom-rank}. Stability of entropic OT with respect to \(\log K\), followed by approximate Sinkhorn scaling and rounding, gives the two error bounds. The low-rank factorization costs \(O(NR^2)\) to construct, each kernel product costs \(O(NR)\), and the quantitative scaling bound contributes \(\widetilde O(NRD^4/(\epsilon\tau))\), yielding the stated time and memory estimates. Their entropy convention differs from the KL-normalized value \(\MK_c^\epsilon\) only by a constant depending on \((a,b)\), so the approximation bounds are unchanged.
\end{proof}
\paragraph{Positive features and complete positivity.}
\begin{defn}[Doubly nonnegative and completely positive matrices and kernels]\label{def-dnn-cp-kernels}
In finite dimension, set
\[
\mathrm{DNN}_n=\enscond{A\in\RR^{n\times n}}{A\succeq0,\ A_{i,j}\geq0},
\qquad
\mathrm{CP}_n=\enscond{BB^\top}{B\in\RR_+^{n\times q}\hbox{ for some }q\geq1}.
\]
Matrices in these sets are called doubly nonnegative and completely positive, respectively. A kernel is doubly positive, respectively completely positive, if every finite Gram matrix belongs to \(\mathrm{DNN}_n\), respectively \(\mathrm{CP}_n\).
It is a positive-feature kernel if there are a probability space \((Z,\xi)\), an integer \(p\geq1\), and measurable features \(\phi(x,\cdot):Z\to\RR_+^p\) such that
\(k(x,y)=\int_Z\dotp{\phi(x,z)}{\phi(y,z)}_{\RR^p}\d\xi(z) \qquad (x,y\in\Xx).\)
\end{defn}
\begin{prop}[Complete positivity and positive features]
\label{prop-cp-kernel-positive-features}
Let \(k:\Xx\times\Xx\to\RR\) be a finite-valued PSD kernel, and assume that its canonical pseudometric
\(d_k(x,y)^2:=k(x,x)+k(y,y)-2k(x,y)\)
is separable. Then \(k\) is completely positive if and only if it is a positive-feature kernel. Moreover, scalar features \(p=1\) suffice. For a finite set \(\Xx=\{x_1,\ldots,x_n\}\), this reduces to
\(K\in\mathrm{CP}_n \quad\Longleftrightarrow\quad K=BB^\top\quad\text{for some }B\geq0.\)
\end{prop}
\begin{proof}
Suppose first that \(k\) has a positive-feature representation. For \(x_1,\ldots,x_n\), let \(v_q(z):=(\phi_q(x_i,z))_{i=1}^n\geq0\). Their Gram matrix is \(K=\sum_{q=1}^p\int_Z v_q(z)v_q(z)^\top\d\xi(z).\)
The integrand belongs to \(\mathrm{CP}_n\), and the closed convex cone \(\mathrm{CP}_n\) is stable under such integrals; hence \(K\in\mathrm{CP}_n\).

Conversely, choose a \(d_k\)-dense sequence \((x_i)_{i\geq1}\) and positive weights \((\omega_i)_{i\geq1}\) such that \(S:=\sum_i\omega_i k(x_i,x_i)<+\infty\). The case \(S=0\) is trivial. For each sufficiently large \(r\), complete positivity gives a factorization \(K^{(r)}=B^{(r)}(B^{(r)})^\top\) of the Gram matrix on \(x_1,\ldots,x_r\). Write \(b_\ell\) for its nonzero columns and set
\(t_\ell:=\sum_{i=1}^r\omega_i(b_\ell)_i^2, \qquad S_r:=\sum_{i=1}^r\omega_i k(x_i,x_i)=\sum_\ell t_\ell.\)
On the compact metrizable product \(Z:=\prod_{i\geq1}[0,\sqrt{S/\omega_i}]\), define \(z^{(r,\ell)}_i=\sqrt{S_r/t_\ell}(b_\ell)_i\) for \(i\leq r\) and \(z^{(r,\ell)}_i=0\) otherwise. Then
\(\xi_r:=\sum_\ell\frac{t_\ell}{S_r}\delta_{z^{(r,\ell)}}\)
is a probability measure on \(Z\), and \(\int_Zz_i z_j\d\xi_r(z)=k(x_i,x_j)\) whenever \(i,j\leq r\). Compactness yields a weakly convergent subsequence with limit \(\xi\). Since every coordinate product \(z_i z_j\) is continuous and bounded on \(Z\), the scalar features \(\phi(x_i,z):=z_i\) satisfy
\(\int_Z\phi(x_i,z)\phi(x_j,z)\d\xi(z)=k(x_i,x_j) \qquad (i,j\geq1).\)
Finally, \(\|\phi(x_i,\cdot)-\phi(x_j,\cdot)\|_{L^2(\xi)}=d_k(x_i,x_j)\). The feature map therefore extends by continuity to all of \(\Xx\); its values remain nonnegative because \(L^2_+(\xi)\) is closed. The Cauchy--Schwarz inequality for PSD kernels shows that the extended inner products equal \(k(x,y)\), which proves the positive-feature representation.
\end{proof}
\(k_\lambda(x,y)=\lambda+\cos^2\pa{\frac{x-y}{2}}, \qquad \lambda>0.\)
\(\langle H,K\rangle = \frac{21}{4}-\frac{5\sqrt5}{2}<0.\)
For a rectangular source--target matrix, Proposition~\ref{prop-cp-kernel-positive-features} is applied to the Gram matrix on the union of the two sampled point sets, after which one extracts the rectangular cross block.

Despite this certification difficulty, completely positive kernels, equivalently positive-sketchable kernels under the assumptions of Proposition~\ref{prop-cp-kernel-positive-features}, enjoy a useful closure algebra. It provides a systematic way to construct new kernels whose low-rank features remain nonnegative.

\begin{prop}[Basic closure of completely positive kernels]
\label{prop-completely-positive-kernel-closure}
Let \(k_1,\ldots,k_J\) be completely positive kernels, \(a_1,\ldots,a_J\geq0\), and \((k_\theta)_{\theta\in\Theta}\) a measurable family of completely positive kernels. Whenever they are finite, the kernels
\[
k_{\mathrm{sum}}:=\sum_{j=1}^J a_jk_j,
\qquad
k_{\mathrm{prod}}:=\prod_{j=1}^J k_j,
\qquad
k_{\mathrm{mix}}:=\int_\Theta k_\theta\,\d\eta(\theta),
\quad \eta\geq0,
\]
are completely positive. In particular, so is every positive autocorrelation
\(k(x,y)=\int g(u-x)g(u-y)\,du, \qquad g\geq0,\)
and every nonnegative mixture of such kernels.
\end{prop}
\begin{proof}
Restrict the kernels to arbitrary points \(x_1,\ldots,x_n\). If \(K_1=B_1B_1^\top\) and \(K_2=B_2B_2^\top\) with \(B_1,B_2\geq0\), then a nonnegative weighted sum is factored by concatenating the correspondingly rescaled columns. Moreover,
\(K_1\odot K_2 = \sum_{\ell,s} \big((B_1)_{\cdot,\ell}\odot(B_2)_{\cdot,s}\big) \big((B_1)_{\cdot,\ell}\odot(B_2)_{\cdot,s}\big)^\top,\)
so the pointwise product is completely positive. A measurable nonnegative mixture is an integral of matrices in the closed convex cone \(\mathrm{CP}_n\), and therefore remains in that cone. Finally, the autocorrelation formula is itself a scalar positive-feature representation, so its finite Gram matrices are completely positive by the same integral argument.
\end{proof}
\begin{prop}[Generalized Gaussian positive features]
\label{prop-gaussian-positive-features}
Let \(0<p\leq2\), \(\epsilon>0\), and
\(k_{p,\epsilon}(x,y) = \exp\pa{-\frac{\norm{x-y}^p}{\epsilon}}, \qquad x,y\in\RR^d.\)
Set \(a=p/2\), let \(S_a\) be the positive \(a\)-stable random variable normalized by \(\EE\pa{e^{-tS_a}}=e^{-t^a}, \qquad t\geq0,\)
and draw independently
\(\Lambda=\epsilon^{-1/a}S_a, \qquad Z\sim\Gaussian(0,\Id_d).\)
For any \(x_0\in\RR^d\), define the positive feature
\(\phi_{p,\epsilon}(x;\Lambda,Z) := \exp\pa{ \sqrt{2\Lambda}\dotp{Z}{x-x_0} -2\Lambda\norm{x-x_0}^2 }.\)
Then \(k_{p,\epsilon}(x,y)=\EE[\phi_{p,\epsilon}(x;\Lambda,Z)\phi_{p,\epsilon}(y;\Lambda,Z)]\), so \(k_{p,\epsilon}\) is completely positive. For i.i.d.\ copies \((\Lambda_\ell,Z_\ell)_{\ell=1}^r\), define the positive sketching features
\(\varphi_\ell(x)=r^{-1/2}\phi_{p,\epsilon}(x;\Lambda_\ell,Z_\ell), \qquad 1\leq\ell\leq r.\)
They satisfy \(\varphi_\ell\geq0\) and \(\EE[\sum_{\ell=1}^r\varphi_\ell(x)\varphi_\ell(y)]=k_{p,\epsilon}(x,y)\). For \(p=2\), one has \(S_1=1\) almost surely and \(\Lambda=1/\epsilon\), hence
\(\phi_{2,\epsilon}(x;Z) = \exp\pa{ \sqrt{\frac{2}{\epsilon}}\dotp{Z}{x-x_0} -\frac{2}{\epsilon}\norm{x-x_0}^2 }.\)
\end{prop}
\begin{proof}
The stable-law normalization gives \(\EE e^{-t\Lambda} = \exp\pa{-\frac{t^a}{\epsilon}},\)
while the Gaussian moment-generating function gives, conditionally on \(\Lambda\), \(\EE_Z\pa{ \phi_{p,\epsilon}(x;\Lambda,Z) \phi_{p,\epsilon}(y;\Lambda,Z) } = e^{-\Lambda\norm{x-y}^2}.\)
Taking \(t=\norm{x-y}^2\) in the first identity proves the feature formula. The Monte Carlo statement follows by averaging independent copies.
\end{proof}
\paragraph{Positive sketches for Sinkhorn.}
\(K_{i,j}=k_{p,\epsilon}(x_i,y_j) =\exp\pa{-\frac{\norm{x_i-y_j}^p}{\epsilon}}.\)
For \(p\geq1\), this is the usual \(p\)-Wasserstein power cost; for \(0<p<1\), it is simply a concave power transport cost. The factor \(1/p\), often used in \(\Wass_p^p\), only rescales \(\epsilon\). Using the sketching features \(\varphi_\ell\) of Proposition~\ref{prop-gaussian-positive-features}, define directly

\((\Phi_X)_{i,\ell}=\varphi_\ell(x_i), \qquad (\Phi_Y)_{j,\ell}=\varphi_\ell(y_j).\)
\(P_r = \diag(u)\Phi_X\Phi_Y^\top\diag(v) = LR^\top, \qquad L=\diag(u)\Phi_X, \quad R=\diag(v)\Phi_Y.\)
\paragraph{Connection with linear time attention.}